\documentclass[10pt,a4paper]{amsart}
\usepackage[margin=19mm]{geometry}
\usepackage{amsmath,amssymb,mathtools}
\usepackage[scr=rsfs,cal=euler]{mathalfa}
\usepackage{xfrac}
\usepackage[unicode,hidelinks,bookmarks=false]{hyperref}
\newcommand{\ie}{i.e.\ }
\newcommand{\cf}{cf.\ }
\newcommand{\resp}{respectively\ }
\newcommand{\Subsection}{\subsection}
\providecommand{\nobreakdash}{\nobreak\hbox{-}}
\setlength{\emergencystretch}{2em}
\multlinegap0pt
\usepackage{amssymb}
\usepackage{enumerate}
\usepackage{tikz}
\usepackage{float}
\usepackage{csquotes}
\usepackage[shortlabels]{enumitem}
\newtheorem{theorem}{Theorem}
\newcommand{\OP}{\mathcal{O}}
\newcommand{\Ptau}{P_\tau}
\title{Order and chain polytopes of \\ maximal ranked posets}
\author{Ibrahim Ahmad, Ghislain Fourier, Michael Joswig}
\date{Source: arXiv:2309.01626v2 (CC BY 4.0)}
\begin{document}
\maketitle

\noindent\emph{Source and licence.} Order and chain polytopes of \\ maximal ranked posets. Version arXiv:2309.01626v2 (CC BY 4.0), distributed by arXiv under the Creative Commons Attribution 4.0 International licence, http://creativecommons.org/licenses/by/4.0/, as recorded in the arXiv licence metadata for this exact version and shown on the abstract page https://arxiv.org/abs/2309.01626v2. Full licence notice: \texttt{licenses/arxiv-2309.01626-CC-BY-4.0.txt}.\par\smallskip

\noindent\textit{Editorial scope.} Counted unit: the article's theorem Theorem-Order-Face-Partition (source0.tex lines 260--269). The article's complete original proof of it is delivered with it (source0.tex lines 389--418, one \texttt{proof} environment of 30 lines, which the article places away from the statement). No mathematics is written, repaired or re-derived: the statement and the proof are the article's own lines, in the article's own notation, and no retained line is substituted or re-flowed.  No further source-local context is needed: every cross-reference of every retained line resolves inside the two delivered spans.  Nothing was omitted from any retained span, so the package carries no omission record.  No retained line contains a citation, so the wrapper carries no bibliography and no bibliography entry.  No retained line contains a figure environment, an image-inclusion command or drawing code, so no image is delivered and the package declares no asset dependency.  Editorial formatting only, in addition: an AMS \texttt{amsart} preamble that restates the article's own declarations and macros used by the retained lines, namely the environments \texttt{theorem} on separate counters, together with the standard packages those lines need.  The retained environments print in this document as Theorem 1 in order of appearance, whereas the article prints the same items with the numbering of its own full document.  The complete native source of the article as distributed in the arXiv e-print for this exact version is bundled unchanged as \texttt{sources/146510/source0.tex}.
\begin{theorem}\label{Theorem-Order-Face-Partition}
Let $P$ be a poset. 
Then every non-empty face $F\subseteq\OP(P)$ can be identified with a partition $\pi_F$ of $\hat{P}$ that is connected, $\hat{P}$-compatible such that $\hat{0}$ and $\hat{1}$ are in different blocks.
Further, the map
\[
  \OP(\hat{P}/\pi_F)\rightarrow F \,,\ (x_{B})_{B\in (\hat{P}/\pi_{F})}\mapsto(x_{q(p)})_{p\in \hat{P}}
\]
is an affine isomorphism.
In particular, we have $\dim(F)=|\pi_F|-2$.
\end{theorem}

\begin{proof}
Checking that $\pi$ is a face partition will be left out as it is a straight-forward verification.
Then according to Theorem \ref{Theorem-Order-Face-Partition}, using the projection onto the blocks $q:P_\tau\rightarrow P_\tau/\pi$, we get the affine isomorphism
\[
   \OP_{C,\varphi}(P_\tau/\pi)\rightarrow F \,,\ (x_{p'})_{p'\in (\Ptau/\pi)}\mapsto (x_{q(p)})_{p\in \Ptau} \enspace .
\]

The dimension of $F$ is exactly $|\pi|-2$, where we account for the blocks containing $\hat{0}$ and $\hat{1}$.
% However, $\hat{0}$ and all elements of $C$ are in singleton blocks.
% For streamlining
% 
However, $\hat{0}$ and all other elements of $C$ are in singleton blocks.
Hence
% \begin{equation*}
%     |\pi|-2=|C|+|\varphi|-1
% \end{equation*}
% For streamlining
% 
\begin{equation*}
    |\pi|-2=|C|+|\varphi|-2
\end{equation*}
% where we account for the block in $\varphi$ containing $\hat{1}$.
% For streamlining
% 
where we account for the block in $\varphi$ containing $\hat{1}$ and $\hat{0}\in C$.
Since the chain-order polytope is of dimension $|O|+|C|-2$, the codimension equals
\begin{equation*}
    |O|+|C|-2-(|C|+|\varphi|-2)=|O|-|\varphi| \enspace . \qedhere
\end{equation*}
\end{proof}

\end{document}
